% introduction.tex -- Introduction

Deep learning has transformed single-chain protein structure prediction, most visibly through AlphaFold2~\cite{jumperHighlyAccurateProtein2021} and RoseTTAFold~\cite{baekRoseTTAFold2021}. The frontier has since shifted toward \emph{designing} novel protein sequences with desired properties~\cite{kuhlmanBradleyAdvances2019,huangComingAgeDeNovo2016}. Generative approaches to this problem now range from classical profile hidden Markov models (pHMMs)~\cite{hmmer3,kroghHiddenMarkovModels1994} and Potts models~\cite{morcosDirectCouplingAnalysis2011,ekeberg2013improved,trinquierEfficientGenerativeModeling2021} to deep generative models including variational autoencoders~\cite{kingmaWellingAutoEncoding2014,riesselman2018deep}, autoregressive language models~\cite{ferruzProtGPT2DeepUnsupervised2022,madaniProGenLanguageModeling2023}, alignment-processing transformers~\cite{raoMSATransformer2021}, MSA-conditioned diffusion~\cite{alamdariProteinGeneration2023}, and structure-conditioned methods such as ProteinMPNN~\cite{dauparasProteinMPNN2022} and RFdiffusion~\cite{watsonRFdiffusion2023}. Yet all sequence-level approaches share a common requirement: large training corpora. Language models are pretrained on millions of UniRef50 sequences; even MSA-conditioned methods require either extensive pretraining or large input alignments. This creates a gap for small families. The median Pfam seed alignment, the curated set of representative sequences that alignment-based methods take as input, contains only 22 sequences. Nearly half of all families have fewer than 20 sequences (Appendix~\ref{si:pfam-census}), below the minimum viable training set for any deep generative architecture. In this scarce-data regime, a few dozen sequences cannot constrain the thousands to millions of parameters in a deep generative model, while pHMMs can emit sequences with low identity to the source family. What is needed is an approach that can extract the statistical structure of a protein family from whatever data is available, even a few dozen sequences, without learning parameters that could lead to overfitting.

The connection between attention and associative memory supplies this approach. Ramsauer et al.~\cite{ramsauerHopfieldNetworksAll2021} showed that a single softmax attention operation is equivalent to one step of the retrieval dynamics on the modern Hopfield energy, building on the dense associative memory framework of Krotov and Hopfield~\cite{krotovDenseAssociativeMemory2016}. This connection endows any set of stored examples with a well-defined energy landscape: the examples become memory patterns, and the Boltzmann distribution over this landscape assigns high probability to states that resemble the stored set. The energy itself has a simple structure. A quadratic term pulls a candidate state toward the origin, while a competing log-sum-exp term over its similarities to the stored patterns pulls it toward whichever memories it most resembles. An inverse temperature controls how sharply the nearest pattern dominates. The gradient of the log-sum-exp is precisely the softmax-weighted combination of stored patterns, i.e., a single attention operation. The score function of the Boltzmann distribution, the gradient that guides sampling toward higher-probability states, is therefore the \emph{residual} of attention: the difference between the attention output and the current state. Diffusion-based generative models such as RFdiffusion~\cite{watsonRFdiffusion2023} obtain such a gradient by training a neural network on data; here it is instead exact and available in closed form, so Langevin sampling can proceed without learning any parameters or approximating gradients. We build on this result to propose \emph{stochastic attention} (SA), a training-free protein sequence generator that combines the exact Hopfield score function with Langevin dynamics. Our prior work~\cite{alswaidanVarnerStochasticAttention2026} introduced SA and validated it on synthetic data, MNIST, and a single protein family. The present study extends that proof of concept from a single illustrative family to a systematic evaluation across diverse families, adding the scaling analysis, baseline benchmarking, and biological validation that the earlier work did not attempt.

In this study, we explored whether stochastic attention can serve as a practical generative model for small protein families. We systematically evaluated SA across eight diverse Pfam families spanning a ten-fold range of family sizes and a seven-fold range of sequence lengths, assessing compositional fidelity, sequence novelty, and structural plausibility using both ESMFold and AlphaFold2. We characterized how generation quality scaled with family size, uncovering an empirical relationship that predicts the critical temperature from alignment statistics alone and eliminates the need for a temperature sweep. We tested robustness in a scarce-data regime down to $K{=}20$ sequences per family and compared SA against four established baselines (profile HMMs, EvoDiff, MSA Transformer, and a Potts model). We also examined whether the generated sequences preserved the pairwise residue covariation that reflects structural contacts and functional couplings. Finally, we sought independent validation of biological plausibility through ESM2-650M pseudo-perplexity scoring and cross-referencing against published deep mutational scanning datasets. Across all eight families, SA generated novel sequences that preserved family-level amino acid statistics and were predicted by both ESMFold and AlphaFold2 to adopt the target-family fold, a combination not matched by the learned baselines tested here. Together these results indicate that the statistical structure of a protein family can itself support sequence generation, without iterative parameter optimization, backpropagation, or external pretraining, when the sampler respects the geometry of the family's sequence space.
